\chapter{Elastography}
\section*{What Is This Test?} Elastography measures tissue stiffness, most commonly in the liver, to estimate fibrosis and sometimes inflammation.
\section*{Alternative Names} Liver elastography; transient elastography; shear-wave elastography; FibroScan.
\section*{Principle} Mechanical waves travel at different speeds through soft and stiff tissue; ultrasound or MRI calculates stiffness.
\section*{Why This Test Is Ordered} It stages chronic liver disease, helps assess fatty liver or viral hepatitis, and may reduce the need for biopsy in selected patients.
\section*{Primary vs Secondary Test} It is a noninvasive secondary assessment; liver biopsy remains useful when diagnosis or staging is uncertain or another disease is suspected.
\section*{How This Test Is Performed} An ultrasound probe is placed on the skin, or MRI sequences measure the liver. The patient holds still and may briefly hold their breath.
\section*{What Preparation Is Needed} Fasting is often requested. Follow local instructions and report ascites, obesity, or implanted devices for MRI-based testing.
\section*{Understanding Results} Higher stiffness can indicate fibrosis but is also affected by inflammation, congestion, cholestasis, and recent meals. Cutoffs vary by method and disease.
\section*{Follow-up Tests} Liver enzymes, viral tests, ultrasound, MRI, biopsy, or hepatology consultation may follow.
\section*{References} \begin{enumerate}\item MedlinePlus. Elastography.\item American Association for the Study of Liver Diseases. Liver fibrosis assessment guidance.\end{enumerate}
\clearpage\section*{Understanding Results and Follow-up}
The laboratory report should identify the specimen, method, units, reference interval or decision limit, and whether the result is positive, negative, abnormal, or inconclusive. Reference intervals vary with age, sex, pregnancy, specimen type, assay, and laboratory; a result outside a range is not automatically a diagnosis, and a result within a range does not always exclude disease. Discuss unexpected results with the ordering clinician, who can decide whether repeat or confirmatory testing, treatment, or referral is appropriate. Do not change medicines or delay urgent care based on an isolated result.
\section*{Regulatory Status and Authorization Date}This chapter describes a test category, not a specific commercial product. FDA, CDSCO, EU IVDR, and MHRA approval or registration dates therefore cannot be assigned without the assay manufacturer, product name, instrument, and intended use. For laboratory-developed tests, an individual agency approval date may not exist. See the Preface for the interpretation of regulatory status.
\clearpage
